\documentclass[11pt]{article}
\usepackage[utf8]{inputenc}
\usepackage[T1]{fontenc}
\usepackage{geometry}
\usepackage{hyperref}
\usepackage{booktabs}
\usepackage{amsmath}
\geometry{margin=1in}
\title{Long Atelier Shell Soak Article}
\author{Atelier Test Suite}
\date{\today}

\begin{document}
\maketitle

\begin{abstract}
This deliberately long document exercises scrolling, wrapping, compilation,
outline navigation, comments, error markers, and synchronization between the
source editor and the generated PDF. The prose is repetitive by design so that
the document remains safe to edit during an extended user soak.
\end{abstract}

\tableofcontents
\newpage

\section{Introduction}
Long-form editing behaves differently from a tiny fixture. Cursor movement,
viewport restoration, line wrapping, and active-section tracking must remain
stable after hundreds of lines have been visited. This article therefore
contains many sections with paragraphs long enough to wrap several times.

The experimental shell should preserve the editor width while the PDF preview
is resized. Closing and reopening the document should restore preferences
without coupling volatile state to another open tab.

\subsection{Motivation}
A realistic soak needs enough material to reveal delayed layout defects. Short
examples rarely expose stale measurements, misplaced gutters, or synchronization
errors that only appear after scrolling far from the beginning of a document.

\subsection{Objectives}
The primary objective is to test ordinary work: writing text, saving changes,
compiling, reading diagnostics, and moving between source and PDF. A secondary
objective is to confirm that comments remain attached after paragraph reflow.

\section{Background}
Document editors combine several independent systems. The source surface owns
text and selections, the compiler reads files from disk, the PDF surface renders
pages, and SyncTeX maps positions between those representations.

Each system may finish asynchronously. A robust interface must prevent a late
response from repainting a destroyed tab or replacing a more important status
message with unrelated background information.

\subsection{Source model}
The source model stores the current text, cursor, selections, and undo history.
Edits should mark the document as modified immediately. Compilation should save
the current buffer before asking the backend to read the file from disk.

\subsection{Preview model}
The preview model may be refreshed after every successful compilation. Its
width is independent from the source model and is controlled by the draggable
separator in the unified shell.

\section{Methods}
The soak procedure begins by opening this file from the gallery with the
experimental LaTeX route enabled. The user then resizes the split, edits several
paragraphs, saves, compiles, and navigates through the outline.

\subsection{Editing protocol}
First, place the cursor in this paragraph and insert a short sentence. Save the
document with the keyboard shortcut and confirm that the modified indicator
returns to the saved state. Undo and redo the edit to exercise history.

Second, select a phrase and add an anchored comment. Rewrap the paragraph at a
different column width, reload the tab, and verify that the comment remains on
the intended phrase.

\subsection{Compilation protocol}
Compile while the buffer contains an unsaved edit. The shell must save before
compilation, and the PDF should contain the new sentence. Repeated clicks on the
compile control must start only one save-and-build chain.

\subsection{Synchronization protocol}
Use forward synchronization from a line near the middle of the document. Then
use reverse synchronization from the PDF and confirm that the editor returns to
the corresponding source line without losing the current split width.

\section{Results}
The following subsections provide enough vertical distance to test navigation
and scrolling. They contain stable labels that can also be used for comments.

\subsection{Layout stability}
The unique phrase LAYOUT-STABILITY-ANCHOR belongs to this subsection. Resize the
separator repeatedly and verify that neither pane collapses below a usable
width. The permitted range should keep both panes visible.

\subsection{Status stability}
The unique phrase STATUS-STABILITY-ANCHOR belongs here. A successful compile
status should remain primary even if background version persistence reports a
secondary warning.

\subsection{Comment stability}
The unique phrase COMMENT-STABILITY-ANCHOR is intended for an anchored comment.
After rewrapping this paragraph, the annotation should still identify the same
text rather than an obsolete line and column pair.

\subsection{Multiple tabs}
Open this article beside the long report fixture. Change the wrap preference in
one tab and move the cursor in the other. Volatile editor state must remain
independent while intentionally shared preferences remain synchronized.

\section{Extended Discussion}
This section adds realistic length. An editor used for manuscripts must remain
responsive after many source lines, repeated compiles, and several preview
refreshes. The test is not a benchmark, but obvious pauses or layout jumps are
signals worth recording during the soak.

Paragraph one discusses navigation. Moving through the outline should place the
cursor at the selected heading and focus the source editor. The outline should
then close so it does not obscure subsequent edits.

Paragraph two discusses diagnostics. A broken command introduced temporarily
below should create a visible error marker and a clickable log entry. Remove the
command afterward and compile again to ensure old markers are cleared.

Paragraph three discusses teardown. Start a compilation and immediately close
the tab. Reopening the file should not reveal stale logs, duplicate listeners,
or controls that react more than once to a single click.

Paragraph four discusses persistence. The split percentage is a user preference
and should survive a reload. Hiding the preview should expand the editor, while
showing it again should restore the last useful percentage.

Paragraph five discusses accessibility. The separator and controls need stable
labels so that keyboard and assistive-technology users can understand their
purpose even when icons are used in the visible interface.

\section{A Small Quantitative Example}
The following equation gives the editor several mathematical lines to render:
\begin{equation}
  \bar{x} = \frac{1}{n}\sum_{i=1}^{n}x_i,
  \qquad
  s^2 = \frac{1}{n-1}\sum_{i=1}^{n}(x_i-\bar{x})^2.
\end{equation}

For a simple linear model,
\begin{equation}
  y_i = \beta_0 + \beta_1 x_i + \varepsilon_i,
  \qquad \varepsilon_i \sim \mathcal{N}(0,\sigma^2).
\end{equation}

\begin{table}[h]
\centering
\begin{tabular}{lrr}
\toprule
Scenario & Source lines & Expected pages \\
\midrule
Short fixture & 20 & 1 \\
Long article & 150+ & 4+ \\
Long report & 180+ & 5+ \\
\bottomrule
\end{tabular}
\caption{Illustrative soak documents.}
\end{table}

\section{Conclusion}
This long article is successful when ordinary editing remains predictable over
an extended session. Record any incorrect status, misplaced annotation, broken
split, stale PDF, or navigation mismatch before enabling the shell by default.

\appendix
\section{Checklist}
\begin{enumerate}
  \item Resize the editor and PDF panes.
  \item Hide and restore the split preview.
  \item Edit, save, undo, and redo.
  \item Compile a dirty buffer.
  \item Navigate through the outline.
  \item Add a comment and rewrap its paragraph.
  \item Exercise forward and reverse SyncTeX.
  \item Open a second LaTeX tab and compare behavior.
  \item Close a tab during compilation.
  \item Reload and verify persisted preferences.
\end{enumerate}

\end{document}
